\chapter{Complexe getallen}\label{ch:b1:complex}

In het bovenbouwvolume kwamen de complexe getallen langs als rekenmiddel
voor vierkantsvergelijkingen. Dit hoofdstuk behandelt ze als een centraal
object: de exponentiële vorm en haar gevolgen (de Moivre, $n$-de wortels,
\omterm{def:b1:complex:unity}{eenheidswortels}), de systematische vertaling tussen $\C$ en de vlakke
meetkunde, en het gebruik van $\eu^{\iu\theta}$ als machine om
goniometrische identiteiten te bewijzen.

\section{Het lichaam $\C$, modulus en toegevoegde}

\begin{definition}[Het complexe lichaam]\label{def:b1:complex:field}
$\C = \{a + \iu b : a, b \in \R\}$ met de gewone optelling en de
vermenigvuldiging die door $\iu^2 = -1$ vastligt. Elke $z = a + \iu b$
ongelijk aan nul heeft een inverse: $z^{-1} = \frac{a - \iu b}{a^2 +
b^2}$ --- in de taal van \cref{ch:b1:structures} is $\C$ een lichaam. Men
schrijft $a = \Re(z)$, $b = \Im(z)$, $\conj{z} = a - \iu b$ (de
\emph{toegevoegde}\index{toegevoegde}) en $\abs{z} = \sqrt{a^2 + b^2}$
(de \emph{modulus}\index{modulus}).
\end{definition}

\begin{proposition}[Rekenregels voor toegevoegde en modulus]\label{prop:b1:complex:rules}
Voor $z, w \in \C$ geldt:
\begin{enumerate}
  \item $\conj{z + w} = \conj z + \conj w$, $\conj{zw} = \conj z\,
        \conj w$, $\conj{\conj z} = z$;
  \item $z \conj z = \abs z^2$; $\;\Re(z) = \frac{z + \conj z}{2}$,
        $\Im(z) = \frac{z - \conj z}{2\iu}$;
  \item $\abs{zw} = \abs z\, \abs w$, en $\abs{z^{-1}} = \abs{z}^{-1}$
        voor $z \neq 0$;
  \item (driehoeksongelijkheid\index{driehoeksongelijkheid}) $\abs{z +
        w} \leq \abs z + \abs w$, met gelijkheid dan en slechts dan als
        $z$ en $w$ op eenzelfde halfrechte vanuit $0$ liggen ($w =
        \lambda z$ of $z = \lambda w$ met $\lambda \geq 0$);
  \item (omgekeerde driehoeksongelijkheid) $\bigl|\abs z - \abs
        w\bigr| \leq \abs{z - w}$.
\end{enumerate}
\end{proposition}

\begin{proof}
(1) en (2) zijn rechtstreekse berekeningen op reële en imaginaire delen.
(3): $\abs{zw}^2 = zw\,\conj{zw} = z\conj z\, w \conj w = \abs z^2 \abs
w^2$, waarna je de wortel trekt; pas dit toe op $z \cdot z^{-1} = 1$ voor
de inverse.

(4) Beide leden zijn niet-negatief, dus vergelijk de kwadraten:
\[
\abs{z+w}^2 = (z+w)(\conj z + \conj w)
= \abs z^2 + \abs w^2 + 2\,\Re(z \conj w),
\]
en uit $\Re(z\conj w) \leq \abs{z \conj w} = \abs z \abs w$ volgt
$\abs{z+w}^2 \leq (\abs z + \abs w)^2$. Gelijkheid dwingt $\Re(z \conj w)
= \abs{z \conj w}$ af, oftewel $z \conj w \in \R_+$; is $w \neq 0$, dan
geeft dat $z = \frac{z\conj w}{\abs w^2}\, w = \lambda w$ met $\lambda
\geq 0$ (en het geval $w = 0$ is triviaal).

(5) Uit $\abs z = \abs{(z - w) + w} \leq \abs{z-w} + \abs w$ volgt $\abs
z - \abs w \leq \abs{z - w}$; verwissel $z$ en $w$ voor het andere teken.
\end{proof}

\begin{example}[Een volledige oefening met modulus en argument]\label{ex:b1:complex:workout}
Breng $w = \dfrac{3 + 4\iu}{1 - 2\iu}$ in algebraïsche vorm en bereken
haar \omterm{def:b1:complex:field}{modulus} tweemaal. Vermenigvuldigen met de \omterm{def:b1:complex:field}{toegevoegde} van de
noemer geeft
\[
w = \frac{(3 + 4\iu)(1 + 2\iu)}{(1 - 2\iu)(1 + 2\iu)}
= \frac{3 + 6\iu + 4\iu - 8}{1 + 4}
= \frac{-5 + 10\iu}{5} = -1 + 2\iu .
\]
Rechtstreeks: $\abs w = \sqrt{1 + 4} = \sqrt5$. Via de quotiëntregel
(\cref{prop:b1:complex:rules}~(3)): $\abs w = \frac{\abs{3 +
4\iu}}{\abs{1 - 2\iu}} = \frac5{\sqrt5} = \sqrt5$ --- hetzelfde
antwoord, zonder algebraïsche vorm. De les veralgemeent: moduli en
argumenten reizen goed door producten en quotiënten, reële en
imaginaire delen reizen goed door sommen. Kies de voorstelling die bij
de voorliggende bewerkingen past, en zet pas om wanneer het niet
anders kan.
\end{example}

\begin{example}[Vergelijkingen met de toegevoegde]\label{ex:b1:complex:conjeq}
Los op in $\C$: $\;z + 2\conj z = 6 + 2\iu$. Een vergelijking waarin
$z$ en $\conj z$ door elkaar lopen is \emph{geen} veeltermvergelijking
in $z$; de betrouwbare zet is opsplitsen in reële coördinaten. Met $z
= x + \iu y$ is
\[
z + 2\conj z = 3x - \iu y ,
\]
zodat de vergelijking $3x = 6$ en $-y = 2$ zegt: de unieke oplossing is
$z = 2 - 2\iu$. (Controle: $(2 - 2\iu) + 2(2 + 2\iu) = 6 + 2\iu$.) Je
kunt ook de hele vergelijking toevoegen, wat $\conj z + 2z = 6 - 2\iu$
oplevert, en het lineaire stelsel in de onbekenden $z, \conj z$
oplossen --- hetzelfde antwoord, en een handige truc zodra de
coëfficiënten complex zijn. Vergelijkingen in $z$ en $\conj z$ zijn in
werkelijkheid stelsels van twee reële vergelijkingen; hier mag je
gerust op ``graad $1$, één oplossing'' rekenen, maar $z\conj z = -1$
(geen oplossing) laat zien dat de veeltermintuïtie het begeeft zodra er
producten opduiken.
\end{example}

\section{Exponentiële vorm}

\begin{definition}[Complexe exponentiële van een imaginair argument]\label{def:b1:complex:expi}
Voor $\theta \in \R$ definieert men
\[
\eu^{\iu\theta} = \cos\theta + \iu \sin\theta .
\]
Elke $z \neq 0$ laat zich schrijven als $z = r\,\eu^{\iu\theta}$ met $r =
\abs z > 0$; zo'n $\theta$ heet een \emph{argument}\index{argument} van
$z$ en ligt vast op een veelvoud van $2\pi$ na. De waarde in
$\intoc{-\pi}{\pi}$ heet het \emph{hoofdargument}, genoteerd $\arg z$.
\end{definition}

\begin{example}[Eerste waarden, en één beroemde identiteit]\label{ex:b1:complex:eulerid}
De definitie gelezen in de hoofdhoeken:
\[
\eu^{\iu\pi/2} = \iu, \qquad
\eu^{\iu\pi} = -1, \qquad
\eu^{2\iu\pi} = 1, \qquad
\eu^{\iu\pi/4} = \frac{\sqrt2}2\,(1 + \iu) .
\]
De tweede, herschreven als $\eu^{\iu\pi} + 1 = 0$, is Eulers beroemde
identiteit die $\eu$, $\iu$, $\pi$, $1$ en $0$ verbindt; in dit
stadium van het boek is ze eerder een definitie die zich ontrolt dan
een stelling, en haar eigenlijke inhoud --- waarom de analytische
exponentiële functie van \cref{ch:b1:functions}, uitgebreid tot
imaginaire argumenten, dezelfde naam verdient --- wordt beslecht door
de machtreeksen van \cref{ch:b1:series}. Ondertussen is bovenstaande
regel het onthouden waard als omrekentabel: ze wordt stilzwijgend
gebruikt telkens als een argument van een tekening wordt afgelezen.
\end{example}

\begin{theorem}[De functionaalvergelijking]\label{thm:b1:complex:funceq}
Voor alle $\theta, \varphi \in \R$ geldt
\[
\eu^{\iu\theta}\, \eu^{\iu\varphi} = \eu^{\iu(\theta + \varphi)},
\qquad
\abs{\eu^{\iu\theta}} = 1,
\qquad
\conj{\eu^{\iu\theta}} = \eu^{-\iu\theta} = (\eu^{\iu\theta})^{-1}.
\]
Bijgevolg draagt $\abs{zw} = \abs z \abs w$ ook de argumenten mee:
$\arg(zw) \equiv \arg z + \arg w \pmod{2\pi}$.
\end{theorem}

\begin{proof}
Werk het product uit en gebruik de somformules:
\[
(\cos\theta + \iu\sin\theta)(\cos\varphi + \iu\sin\varphi)
= (\cos\theta\cos\varphi - \sin\theta\sin\varphi)
+ \iu\,(\sin\theta\cos\varphi + \cos\theta\sin\varphi),
\]
wat gelijk is aan $\cos(\theta+\varphi) + \iu\sin(\theta+\varphi)$. De
\omterm{def:b1:complex:field}{modulus} is $\sqrt{\cos^2\theta + \sin^2\theta} = 1$, en de formule voor
de \omterm{def:b1:complex:field}{toegevoegde} is de pariteit van cosinus en sinus; ze inverteert
$\eu^{\iu\theta}$, want $\eu^{\iu\theta}\eu^{-\iu\theta} = \eu^0 = 1$.
\end{proof}

\begin{corollary}[Formule van de Moivre]\label{cor:b1:complex:demoivre}
Voor $\theta \in \R$ en $n \in \Z$ geldt $\;(\cos\theta +
\iu\sin\theta)^n = \cos(n\theta) + \iu\sin(n\theta)$.\index{formule van
de Moivre}
\end{corollary}

\begin{proof}
Voor $n \geq 0$ met inductie. Het geval $n = 0$ luidt $1 = 1$. Neem de
formule aan voor $n$, dan geeft de functionaalvergelijking
(\cref{thm:b1:complex:funceq})
\[
\bigl(\eu^{\iu\theta}\bigr)^{n+1}
= \bigl(\eu^{\iu\theta}\bigr)^{n}\,\eu^{\iu\theta}
= \eu^{\iu n\theta}\,\eu^{\iu\theta}
= \eu^{\iu(n+1)\theta} ,
\]
wat de formule op rang $n + 1$ is. Voor $n < 0$ schrijf je $n = -m$ met
$m > 0$: omdat $\eu^{\iu m\theta}$ de inverse $\eu^{-\iu m\theta}$ heeft
(opnieuw \cref{thm:b1:complex:funceq}), is
\[
\bigl(\eu^{\iu\theta}\bigr)^{-m}
= \bigl(\eu^{\iu m\theta}\bigr)^{-1}
= \eu^{\iu(-m)\theta} . \qedhere
\]
\end{proof}

\begin{example}[$\cos 3\theta$ uitwerken met de Moivre]\label{ex:b1:complex:cos3}
Schrijf $c = \cos\theta$ en $s = \sin\theta$. De Moivre en het binomium
geven
\[
\cos 3\theta + \iu \sin 3\theta = (c + \iu s)^3
= c^3 - 3cs^2 + \iu\,(3c^2 s - s^3),
\]
en na gelijkstellen van de reële delen en substitutie van $s^2 = 1 -
c^2$:
\[
\cos 3\theta = c^3 - 3c(1 - c^2) = 4\cos^3\theta - 3\cos\theta .
\]
Het imaginaire deel levert er gratis $\sin 3\theta = 3\sin\theta -
4\sin^3\theta$ bij: één complexe identiteit draagt altijd \emph{twee}
reële. Achterstevoren gelezen is de zojuist gevonden identiteit de
sleutel tot de klassieke driedelingsvergelijking: $\cos(\theta/3)$ uit
$\cos\theta$ construeren betekent de derdegraadsvergelijking $4x^3 - 3x
= \cos\theta$ oplossen, en daar neemt de algebra van \cref{ch:b1:poly}
het over.
\end{example}

\begin{proposition}[Formules van Euler]\label{prop:b1:complex:euler}
\[
\cos\theta = \frac{\eu^{\iu\theta} + \eu^{-\iu\theta}}{2},
\qquad
\sin\theta = \frac{\eu^{\iu\theta} - \eu^{-\iu\theta}}{2\iu}.
\]
\end{proposition}

\begin{proof}
Tel $\eu^{\iu\theta} = \cos\theta + \iu\sin\theta$ en $\eu^{-\iu\theta} =
\cos\theta - \iu\sin\theta$ bij elkaar op, respectievelijk trek ze van
elkaar af.
\end{proof}

\begin{method}[Goniometrie via exponentiëlen]\label{met:b1:complex:trig}
\begin{enumerate}
  \item \emph{Lineariseer} $\cos^p\theta \sin^q\theta$ (zet machten om
        in een som van $\cos k\theta$ en $\sin k\theta$): substitueer de
        formules van Euler, werk uit met het binomium
        (\cref{thm:b1:counting:binomial}) en hergroepeer \omterm{def:b1:complex:field}{toegevoegde}
        termen.
  \item \emph{Ontwikkel} $\cos n\theta$ als veelterm in $\cos\theta$:
        schrijf $\cos n\theta = \Re\bigl((\cos\theta +
        \iu\sin\theta)^n\bigr)$, werk uit en zet even machten van
        $\sin$ om met $\sin^2 = 1 - \cos^2$.
  \item \emph{Sommeer goniometrische reeksen} als $\sum_k \cos k\theta$:
        herken het reële deel van een meetkundige som $\sum_k
        (\eu^{\iu\theta})^k$.
  \item De \emph{halvehoekontbinding}: voor alle $p, q$ geldt
        \[
        \eu^{\iu p} + \eu^{\iu q}
        = 2 \cos\tfrac{p - q}{2}\; \eu^{\iu \frac{p+q}{2}},
        \qquad
        \eu^{\iu p} - \eu^{\iu q}
        = 2\iu \sin\tfrac{p - q}{2}\; \eu^{\iu \frac{p+q}{2}} .
        \]
\end{enumerate}
\end{method}

\begin{example}[Lineariseren]\label{ex:b1:complex:linearize}
\[
\cos^3\theta
= \Bigl(\frac{\eu^{\iu\theta} + \eu^{-\iu\theta}}{2}\Bigr)^{\!3}
= \frac{\eu^{3\iu\theta} + 3\eu^{\iu\theta} + 3\eu^{-\iu\theta} +
\eu^{-3\iu\theta}}{8}
= \frac{\cos 3\theta + 3\cos\theta}{4}.
\]
Deze vorm laat zich meteen integreren --- precies de reden waarom
lineariseren ertoe doet in \cref{ch:b1:integration}. Een gemengd
product gaat net zo, alleen met beide formules van Euler tegelijk:
\[
\sin^2\theta\cos^2\theta
= \Bigl(\frac{\sin2\theta}2\Bigr)^{\!2}
= \frac{1}{4}\cdot
\Bigl(\frac{\eu^{2\iu\theta} - \eu^{-2\iu\theta}}{2\iu}
\Bigr)^{\!2}
= \frac{2 - \eu^{4\iu\theta} - \eu^{-4\iu\theta}}{16}
= \frac{1 - \cos4\theta}{8} ,
\]
waarbij de dubbelehoekformule in de eerste stap één binomiale
ontwikkeling uitspaarde --- daar loont het altijd even naar te kijken
vóór je gaat mechaniseren.
\end{example}

\begin{example}[Een binomiale goniometrische som]\label{ex:b1:complex:binomsum}
Bereken voor $n \in \N$ en $\theta \in \R$ de som $S = \sum_{k=0}^{n}
\binom nk \cos k\theta$. Herken het reële deel van een binomiale
ontwikkeling:
\[
S = \Re\sum_{k=0}^n \binom nk \bigl(\eu^{\iu\theta}\bigr)^k
= \Re\bigl(1 + \eu^{\iu\theta}\bigr)^n ,
\]
en ontbind dan de halve hoek (\cref{met:b1:complex:trig}~(4)): $1 +
\eu^{\iu\theta} = 2\cos\frac\theta2\,\eu^{\iu\theta/2}$, zodat
\[
S = \Re\Bigl(2^n\cos^n\frac\theta2\;\eu^{\iu n\theta/2}\Bigr)
= 2^n \cos^n\frac\theta2\,\cos\frac{n\theta}2 .
\]
Het imaginaire deel levert er gratis $\sum_k\binom nk\sin k\theta =
2^n\cos^n\frac\theta2\sin\frac{n\theta}2$ bij. Controles: $\theta = 0$
geeft $\sum\binom nk = 2^n$ terug, en $\theta = \pi$ geeft $S = 0$ voor
$n \geq 1$ (elke factor $\cos\frac\pi2$ verdwijnt), dat wil zeggen de
alternerende rijsom uit \cref{ex:b1:counting:binomial}. De methode ---
``zie de cosinussom als de schaduw van een complexe macht en ontbind
dan halve hoeken'' --- is precies die van \cref{exo:b1:complex:6}, met
het binomium in de plaats van de meetkundige reeks.
\end{example}

\section{Wortels van complexe getallen}

\begin{theorem}[$n$-de wortels]\label{thm:b1:complex:roots}
Zij $a = r\,\eu^{\iu\alpha} \neq 0$ en $n \in \N^*$. De vergelijking
$z^n = a$ heeft precies $n$ oplossingen:
\[
z_k = r^{1/n}\, \eu^{\iu\left(\frac{\alpha}{n} +
\frac{2k\pi}{n}\right)},
\qquad k = 0, 1, \dots, n - 1 .
\]
\end{theorem}

\begin{proof}
Schrijf $z = \rho\,\eu^{\iu\theta}$ met $\rho > 0$. Dan is $z^n = \rho^n
\eu^{\iu n\theta} = r \eu^{\iu\alpha}$ dan en slechts dan als $\rho^n =
r$ (moduli) en $n\theta \equiv \alpha \pmod{2\pi}$ (argumenten), dat wil
zeggen $\rho = r^{1/n}$ en $\theta = \frac{\alpha}{n} + \frac{2k\pi}{n}$
voor zekere $k \in \Z$. Rest de vraag wanneer twee gehele getallen $k,
k'$ hetzelfde getal opleveren: dat gebeurt precies wanneer de hoeken een
veelvoud van $2\pi$ verschillen,
\[
\frac{2k\pi}n - \frac{2k'\pi}n \in 2\pi\Z
\iff \frac{k - k'}n \in \Z
\iff n \mid k - k' .
\]
Wegens de euclidische deling is elke $k \in \Z$ modulo $n$ congruent met
precies één element van $\{0, 1, \dots, n-1\}$, zodat dat bereik elke
oplossing één keer opsomt en de telling exact $n$ oplevert. (Datzelfde
argument, uitgevoerd binnen $\mathbb U_n$, laat zien dat de wortels een
regelmatige $n$-hoek vormen: opeenvolgende waarden van $k$ draaien over
de vaste hoek $\frac{2\pi}n$.)
\end{proof}

\begin{example}[Derdemachtswortels van $-27$]\label{ex:b1:complex:cuberoots}
Los $z^3 = -27$ op. Exponentiële vorm van het rechterlid: $-27 =
27\,\eu^{\iu\pi}$, zodat de drie wortels
\[
z_k = 3\,\eu^{\iu(\frac\pi3 + \frac{2k\pi}3)}, \quad k = 0, 1, 2 :
\qquad
z_0 = 3\eu^{\iu\pi/3} = \frac32 + \frac{3\sqrt3}2\,\iu,
\quad
z_1 = -3,
\quad
z_2 = \conj{z_0}
\]
zijn. Twee controles. Ten eerste is de reële wortel $-3$ de voor de
hand liggende, en de twee andere zijn haar draaiingen over
$\pm\frac{2\pi}3$ --- equivalent: $-3j$ en $-3j^2$. Ten tweede
bevestigt de algebra het: $z^3 + 27 = (z + 3)(z^2 - 3z + 9)$, en de
vierkantsvergelijking heeft discriminant $9 - 36 = -27 < 0$ met
wortels $\frac{3 \pm 3\iu\sqrt3}2 = z_0, \conj{z_0}$. Het inzicht: bij
een reëel rechterlid komen de niet-reële wortels altijd in \omterm{def:b1:complex:field}{toegevoegde}
paren, zodat een tekening van de oplossingsverzameling symmetrisch is
om de reële as --- een voorproefje van de reële ontbindingsstelling
van \cref{ch:b1:poly}.
\end{example}

\begin{example}[Een niet-reëel rechterlid]\label{ex:b1:complex:fourthroots}
Los $z^4 = -8 + 8\iu\sqrt3$ op. Exponentiële vorm van het rechterlid:
\omterm{def:b1:complex:field}{modulus} $\sqrt{64 + 192} = 16$, en argument $\theta$ met $\cos\theta =
-\frac12$ en $\sin\theta = \frac{\sqrt3}2$, dus $\theta =
\frac{2\pi}3$. De vier wortels zijn
\[
z_k = 2\,\eu^{\iu(\frac\pi6 + \frac{k\pi}2)}, \quad k = 0, 1, 2,
3 :
\qquad
z_0 = \sqrt3 + \iu,\quad z_1 = \iu z_0 = -1 + \iu\sqrt3,
\]
\[
z_2 = -z_0 = -\sqrt3 - \iu,\qquad z_3 = -\iu z_0 = 1 - \iu\sqrt3 .
\]
(Controle: $z_0^2 = 2 + 2\iu\sqrt3$, dus $z_0^4 = (2 + 2\iu\sqrt3)^2 =
4 - 12 + 8\iu\sqrt3 = -8 + 8\iu\sqrt3$.) Zodra \emph{één} wortel
gevonden is, komen de drie andere er gratis bij: het zijn haar
opeenvolgende draaiingen over $\frac\pi2$, dat wil zeggen haar
producten met de vierde \omterm{def:b1:complex:unity}{eenheidswortels} --- de algemene structuur
achter \cref{thm:b1:complex:roots}, die je beter benut dan elke wortel
opnieuw uit te rekenen. Ditmaal is er geen \omterm{def:b1:complex:field}{toegevoegde} symmetrie: het
rechterlid is niet reëel.
\end{example}

\begin{definition}[Eenheidswortels]\label{def:b1:complex:unity}
De \emph{$n$-de eenheidswortels}\index{eenheidswortels} zijn de
oplossingen van $z^n = 1$:
\[
\mathbb{U}_n = \bigl\{\, \omega^k : k = 0, \dots, n-1 \,\bigr\},
\qquad \omega = \eu^{2\iu\pi/n}.
\]
Ze vormen een groep onder de vermenigvuldiging
(\cref{ch:b1:structures}) en liggen op de hoekpunten van een
regelmatige $n$-hoek ingeschreven in de eenheidscirkel.
\end{definition}

\begin{omfigure}
\begin{tikzpicture}[scale=1.5]
  \draw[->, thin] (-1.35,0) -- (1.45,0) node[below] {$\Re$};
  \draw[->, thin] (0,-1.35) -- (0,1.45) node[left] {$\Im$};
  \draw[thin, gray] (0,0) circle (1);
  \foreach \k in {0,...,4} {
    \fill[omDef] ({cos(72*\k)},{sin(72*\k)}) circle (0.035);
  }
  \draw[omDef, thick] (1,0) -- ({cos(72)},{sin(72)}) -- ({cos(144)},{sin(144)})
    -- ({cos(216)},{sin(216)}) -- ({cos(288)},{sin(288)}) -- cycle;
  \node[right, font=\small, omDef] at (1.02,0.12) {$1$};
  \node[above right, font=\small, omDef] at ({cos(72)},{sin(72)})
    {$\omega$};
  \node[above left, font=\small, omDef] at ({cos(144)},{sin(144)})
    {$\omega^2$};
  \node[below left, font=\small, omDef] at ({cos(216)},{sin(216)})
    {$\omega^3$};
  \node[below right, font=\small, omDef] at ({cos(288)},{sin(288)})
    {$\omega^4$};
\end{tikzpicture}

{\small De vijfde \omterm{def:b1:complex:unity}{eenheidswortels}, $\omega = \eu^{2\iu\pi/5}$: een
regelmatige vijfhoek op de eenheidscirkel.}
\end{omfigure}

\begin{proposition}[Som van de eenheidswortels]\label{prop:b1:complex:sumroots}
Voor $n \geq 2$ tellen de $n$-de \omterm{def:b1:complex:unity}{eenheidswortels} op tot nul:
$\sum_{k=0}^{n-1} \omega^k = 0$.
\end{proposition}

\begin{proof}
Meetkundige som met reden $\omega \neq 1$: $\sum_{k=0}^{n-1} \omega^k =
\frac{\omega^n - 1}{\omega - 1} = 0$, want $\omega^n = 1$.
\end{proof}

\begin{example}[Reële en imaginaire delen aflezen]\label{ex:b1:complex:sumrootsparts}
Splitsen we $\sum_{k=0}^{n-1}\omega^k = 0$ in reëel en imaginair deel,
dan vallen er gratis twee goniometrische identiteiten uit:
\[
\sum_{k=0}^{n-1}\cos\frac{2k\pi}n = 0,
\qquad
\sum_{k=0}^{n-1}\sin\frac{2k\pi}n = 0
\qquad (n \geq 2).
\]
Meetkundig: het zwaartepunt van een regelmatige $n$-hoek ingeschreven
in de eenheidscirkel is haar middelpunt --- de hoekpunten heffen
elkaar precies op. Voor $n = 5$ geeft de eerste identiteit $1 +
2\cos\frac{2\pi}5 + 2\cos\frac{4\pi}5 = 0$ (door $k$ met $n - k$ te
koppelen), het beginpunt van de berekening van $\cos\frac{2\pi}5$ in
\cref{exo:b1:complex:8}.
\end{example}

\begin{example}[Vierkantswortels in algebraïsche vorm]\label{ex:b1:complex:sqrt}
Om $z^2 = 3 + 4\iu$ zonder goniometrie op te lossen, zet je $z = x + \iu
y$:
\[
x^2 - y^2 = 3, \qquad 2xy = 4, \qquad x^2 + y^2 = \abs{3 + 4\iu} = 5 .
\]
De eerste en de laatste optellen geeft $x^2 = 4$, dus $x = \pm 2$, en dan
$y = 2/x = \pm 1$ met \emph{hetzelfde} tekenpaar (want $xy = 2 > 0$): $z
= \pm(2 + \iu)$. Samen met de gebruikelijke formule lost dit elke
vierkantsvergelijking met complexe coëfficiënten op
(\cref{exo:b1:complex:7}).
\end{example}

\section{Complexe getallen en vlakke meetkunde}

\begin{proposition}[Meetkundig woordenboek]\label{prop:b1:complex:geometry}
Vereenzelvig het punt $M(x, y)$ van het vlak met zijn \emph{affix} $z = x
+ \iu y$\index{affix}. Voor verschillende punten $A, B, C$ met affixen
$a, b, c$:
\begin{enumerate}
  \item $\abs{b - a}$ is de afstand $AB$;
  \item $\arg\dfrac{c - a}{b - a}$ is de hoek tussen de vectoren
        $\vect{AB}$ en $\vect{AC}$ (modulo $2\pi$);
  \item $A, B, C$ liggen op één rechte dan en slechts dan als
        $\dfrac{c - a}{b - a} \in \R$; de rechten $AB$ en $AC$ staan
        loodrecht op elkaar dan en slechts dan als $\dfrac{c - a}{b -
        a} \in \iu\R$.
\end{enumerate}
\end{proposition}

\begin{proof}
(1) is de definitie van de \omterm{def:b1:complex:field}{modulus} toegepast op $b - a$, de affix van
$\vect{AB}$. (2): schrijf $b - a = r\eu^{\iu\theta}$ en $c - a =
s\eu^{\iu\varphi}$; dan heeft $\frac{c-a}{b-a} = \frac sr \eu^{\iu(\varphi
- \theta)}$ als argument $\varphi - \theta$, de hoek van $\vect{AB}$ naar
$\vect{AC}$. (3): op één rechte liggen betekent hoek $0$ of $\pi$, dus
argument in $\pi\Z$, dus quotiënt reëel; loodrecht staan betekent hoek
$\pm\frac\pi2$, dus quotiënt zuiver imaginair. (Het quotiënt is niet nul,
want $C \neq A$.)
\end{proof}

\begin{remark}[Tussenspel: $\C$ is het vlak met een vermenigvuldiging]\label{rem:b1:complex:interlude}
Het loont stil te staan bij wat dit hoofdstuk überhaupt mogelijk
maakt: het vlak $\R^2$ draagt optellingen in elke richting, maar geen
van nature gegeven vermenigvuldiging --- en $\C$ \emph{is} het vlak
uitgerust met er één, waarin vermenigvuldigen met een vast getal
draait en schaalt. Deze ene structuur wordt in dit volume nog driemaal
uitgemolken. In \cref{ch:b1:matrices} duikt vermenigvuldigen met $a +
\iu b$ opnieuw op als de $2 \times 2$-matrix met rijen $(a, -b)$ en
$(b, a)$: het complexe rekenen is een eerste, volstrekt concrete
familie matrixproducten. In \cref{ch:b1:euclid} blijkt de uitdrukking
$\Re(\conj z\,w)$ het inwendig product te zijn en $\abs z$ de
euclidische norm: de hier bewezen driehoeksongelijkheid is het model
voor het verhaal van Cauchy--Schwarz daar. En in \cref{ch:b1:curves}
schrijf je een bewegend punt het best als $t \mapsto z(t)$, zodat
snelheid en versnelling complexwaardige afgeleiden worden --- een
cirkelbeweging is dan simpelweg $z(t) = R\,\eu^{\iu\omega t}$. Eén
goede vermenigvuldiging, vier hoofdstukken rendement.
\end{remark}

\begin{proposition}[De afbeeldingen $z \mapsto az + b$]\label{prop:b1:complex:similitude}
Zij $a \in \C^*$ en $b \in \C$. De transformatie $f(z) = az + b$ van het
vlak:
\begin{itemize}
  \item is een translatie wanneer $a = 1$;
  \item heeft anders precies één vast punt $\zeta = \frac{b}{1-a}$, en
        $f(z) - \zeta = a\,(z - \zeta)$: $f$ is de draaiing met
        middelpunt $\zeta$ en hoek $\arg a$, samengesteld met de
        vermenigvuldiging met middelpunt $\zeta$ en factor $\abs a$.
\end{itemize}
In het bijzonder is $z \mapsto \eu^{\iu\theta} z$ de draaiing over de
hoek $\theta$ om de oorsprong, en is $\conj z$ de spiegeling in de reële
as.
\end{proposition}

\begin{proof}
Is $a = 1$, dan verschuift $f(z) = z + b$ over de vector met affix $b$.
Is $a \neq 1$, dan heeft de vergelijking $z = az + b$ voor de vaste
punten de unieke oplossing $\zeta = \frac{b}{1-a}$, en dan is $f(z) -
\zeta = az + b - (a\zeta + b) = a(z - \zeta)$. Schrijven we $a = \abs a\,
\eu^{\iu\arg a}$, dan schaalt vermenigvuldigen met $a$ de afstanden tot
$\zeta$ met $\abs a$ en telt het $\arg a$ bij de hoeken in $\zeta$ op ---
en dat is de aangekondigde samenstelling.
\end{proof}

\begin{example}[Een afbeelding $z \mapsto az+b$ herkennen]\label{ex:b1:complex:classify}
Neem $f(z) = \iu z + 1$. Hier is $a = \iu \neq 1$, dus het vaste punt is
\[
\zeta = \frac{b}{1 - a} = \frac1{1 - \iu}
= \frac{1 + \iu}{2},
\]
en omdat $\abs a = 1$ wordt er helemaal niet geschaald: $f$ is de zuivere
draaiing met middelpunt $\frac{1+\iu}2$ over de hoek $\arg \iu =
\frac\pi2$. Controle: $f(\zeta) = \iu\,\frac{1+\iu}2 + 1 = \frac{\iu -
1}2 + 1 = \frac{1 + \iu}2 = \zeta$, en $f(0) = 1$, $f(1) = 1 + \iu$, $f(1
+ \iu) = \iu\,(1 + \iu) + 1 = \iu$: de vier punten $0, 1, 1+\iu, \iu$ van
het eenheidsvierkant lopen telkens een kwartslag rond hun middelpunt
$\zeta$ --- precies wat een draaiing over $\frac\pi2$ om het middelpunt
van het vierkant moet doen.
\end{example}

\begin{omfigure}
\begin{tikzpicture}[scale=1.6]
  \draw[->, thin] (-0.5,0) -- (1.6,0) node[below] {$\Re$};
  \draw[->, thin] (0,-0.35) -- (0,1.5) node[left] {$\Im$};
  \draw[gray, thin] (0,0) rectangle (1,1);
  \fill[omDef] (0.5,0.5) circle (0.03)
    node[below right, font=\small] {$\zeta$};
  \fill (0,0) circle (0.025) node[below left, font=\small] {$0$};
  \fill (1,0) circle (0.025) node[below right, font=\small] {$1$};
  \fill (1,1) circle (0.025)
    node[above right, font=\small] {$1{+}\iu$};
  \fill (0,1) circle (0.025) node[above left, font=\small] {$\iu$};
  \draw[omThm, ->, thick] (0.08,-0.12)
    .. controls (0.6,-0.28) and (1.12,0.4) .. (1.12,0.9);
  \draw[omThm, ->, thick] (1.12,1.08)
    .. controls (0.9,1.3) and (0.35,1.3) .. (0.05,1.12);
  \node[omThm, font=\small] at (0.95,-0.28) {$0 \mapsto 1$};
  \node[omThm, font=\small] at (1.62,0.75) {$1 \mapsto 1{+}\iu$};
\end{tikzpicture}

{\small De \omterm{def:b1:logic:map}{afbeelding} $f(z) = \iu z + 1$ uit
\cref{ex:b1:complex:classify}: een kwartslagdraaiing om het vaste punt
$\zeta = \frac{1+\iu}2$. De hoekpunten van het eenheidsvierkant
doorlopen de cyclus $0 \to 1 \to 1{+}\iu \to \iu \to 0$; geen enkel
punt beweegt langs een rechte lijn, en toch wordt het hele vierkant
star gedraaid.}
\end{omfigure}

\begin{remark}[Veelgemaakte fouten met moduli en argumenten]\label{rem:b1:complex:pitfalls}
\begin{enumerate}
  \item \emph{Geen ongelijkheden in $\C$.} $z \leq w$ schrijven voor
        niet-reële getallen is betekenisloos; alleen moduli, reële
        delen en imaginaire delen zijn vergelijkbaar.
  \item \emph{$\sqrt{\phantom z}$ is voorbehouden aan $\R_+$.} Elk
        complex getal ongelijk aan nul heeft \emph{twee}
        vierkantswortels en geen van beide is bevoorrecht: schrijf
        ``zij $\delta$ een vierkantswortel van $\Delta$'' (berekend als
        in \cref{ex:b1:complex:sqrt}), nooit $\sqrt\Delta$ --- de regel
        $\sqrt{ab} = \sqrt a\sqrt b$ faalt al bij $a = b = -1$.
  \item \emph{Argumenten leven modulo $2\pi$.} Uit $\eu^{\iu\alpha} =
        \eu^{\iu\beta}$ besluit je $\alpha \equiv \beta \pmod{2\pi}$ en
        niet $\alpha = \beta$; de $2k\pi$ vergeten is de manier waarop
        oplossingsverzamelingen van $z^n = a$ $n - 1$ van hun $n$
        elementen kwijtraken.
  \item \emph{$\abs{z + w}$ is niet $\abs z + \abs w$.} Gelijkheid in
        de driehoeksongelijkheid is het uitzonderlijke geval waarin
        alles op één halfrechte ligt
        (\cref{prop:b1:complex:rules}~(4)); in het algemeen moet de
        \omterm{def:b1:complex:field}{modulus} van een som afgeschat worden, niet berekend.
\end{enumerate}
\end{remark}

\begin{example}\label{ex:b1:complex:equilateral}
$A, B, C$ (affixen $a, b, c$, twee aan twee verschillend) vormen een
gelijkzijdige driehoek met de hoekpunten in directe (tegen de wijzers
van de klok in) volgorde dan en slechts dan als $\frac{c - a}{b - a} =
\eu^{\iu\pi/3}$: de draaiing met middelpunt $A$ over de hoek $\frac\pi3$
stuurt $B$ naar $C$. Beide oriëntaties samen worden gevat door de
symmetrische vergelijking $a^2 + b^2 + c^2 = ab + bc + ca$
(\cref{exo:b1:complex:10}).
\end{example}

\begin{remark}[Waar dit hoofdstuk gebruikt wordt]\label{rem:b1:complex:whereused}
De exponentiële vorm is het meest hergebruikte rekenmiddel van het hele
volume. \omterm{def:b1:complex:unity}{Eenheidswortels} worden het standaardvoorbeeld van een cyclische
groep in \cref{ch:b1:structures} en dragen de ontbinding van $X^n - 1$
in \cref{ch:b1:poly}; \omterm{def:b1:complex:field}{toegevoegde} wortels daar groeperen levert de reële
ontbindingen die de partieelbreuksplitsing van \cref{ch:b1:fractions}
gebruikt. Lineariseren (\cref{met:b1:complex:trig}) is de
standaardvoorbereiding op het integreren van goniometrische machten in
\cref{ch:b1:integration}, en de karakteristieke vergelijkingen van
\cref{ch:b1:diffeq} hebben complexe wortels waarvan de reële en
imaginaire delen de oscillerende oplossingen $\eu^{\lambda t}\cos\omega
t$ opleveren. Het meetkundige woordenboek keert in matrixgewaad terug:
draaiingen en gelijkvormigheden worden de orthogonale matrices van
\cref{ch:b1:matrices,ch:b1:euclid}, en geparametriseerde krommen in
\cref{ch:b1:curves} schrijf je vaak het best als \omterm{def:b1:logic:map}{afbeeldingen} $t \mapsto
z(t) \in \C$.
\end{remark}

\section{Oefeningen}

\begin{exercise}[$\star$]\label{exo:b1:complex:1}
Breng in exponentiële vorm: $1 + \iu$; $\;\sqrt 3 - \iu$; $\;-5$;
$\;\dfrac{1 + \iu}{\sqrt 3 - \iu}$. Leid daaruit
$\cos\frac{5\pi}{12}$ en $\sin\frac{5\pi}{12}$ af.
\end{exercise}

\begin{exercise}[$\star$]\label{exo:b1:complex:2}
Bereken $(1 + \iu)^{20}$. Voor welke $n \in \N$ is $(1 + \iu)^n$ een
reëel getal?
\end{exercise}

\begin{exercise}[$\star$]\label{exo:b1:complex:3}
Beschrijf meetkundig de \omterm{def:b1:logic:sets}{verzameling} van de $z \in \C$ waarvoor:
$\abs{z - 2} = \abs{z + \iu}$; $\;\abs{z - 1} = 2$;
$\;\dfrac{z - 1}{z + 1} \in \iu\R$ (voor $z \neq -1$).
\end{exercise}

\begin{exercise}[$\star$]\label{exo:b1:complex:4}
Lineariseer $\sin^4\theta$, en ontwikkel $\cos 4\theta$ als veelterm in
$\cos\theta$.
\end{exercise}

\begin{exercise}[$\star$]\label{exo:b1:complex:5}
Los $z^3 = 8\iu$ op en zet de oplossingen op een tekening. Los $z^4 =
-4$ op en ontbind $X^4 + 4$ in twee reële tweedegraadsveeltermen.
\end{exercise}

\begin{exercise}[$\star\star$]\label{exo:b1:complex:6}
Bereken voor $\theta \in \R$ met $\eu^{\iu\theta} \neq 1$ en $n \in \N$
\[
C_n = \sum_{k=0}^{n} \cos k\theta
\qquad\text{en}\qquad
S_n = \sum_{k=0}^{n} \sin k\theta
\]
door de meetkundige reeks $\sum_k \eu^{\iu k\theta}$ te sommeren en de
halvehoekontbinding te gebruiken.
\end{exercise}

\begin{exercise}[$\star\star$]\label{exo:b1:complex:7}
Los op in $\C$: $z^2 - (3 + 4\iu) z + (-1 + 5\iu) = 0$. \emph{(Bereken
de discriminant en trek er de vierkantswortels uit als in
\cref{ex:b1:complex:sqrt}.)}
\end{exercise}

\begin{exercise}[$\star\star$]\label{exo:b1:complex:8}
Zij $\omega = \eu^{2\iu\pi/5}$.
\begin{enumerate}
  \item Verantwoord $1 + \omega + \omega^2 + \omega^3 + \omega^4 = 0$.
  \item Zet $u = \omega + \omega^4$ en $v = \omega^2 + \omega^3$.
        Bereken $u + v$ en $uv$, en leid af dat $u$ en $v$ de wortels
        van $X^2 + X - 1$ zijn.
  \item Besluit dat $\cos\frac{2\pi}{5} = \frac{\sqrt 5 - 1}{4}$.
\end{enumerate}
\end{exercise}

\begin{exercise}[$\star\star$]\label{exo:b1:complex:9}
Bewijs voor alle $z, w \in \C$ de parallellogramidentiteit
\[
\abs{z + w}^2 + \abs{z - w}^2 = 2\abs z^2 + 2\abs w^2 ,
\]
en geef er een meetkundige uitleg bij in het parallellogram met
hoekpunten $0, z, w, z + w$.
\end{exercise}

\begin{exercise}[$\star\star\star$]\label{exo:b1:complex:10}
Bewijs dat drie twee aan twee verschillende punten met affixen $a, b, c$
een gelijkzijdige driehoek vormen (in welke oriëntatie dan ook) dan en
slechts dan als
\[
a^2 + b^2 + c^2 = ab + bc + ca .
\]
\emph{Aanwijzing: het directe geval is $\frac{c-a}{b-a} = -j^2$ en het
indirecte $\frac{c-a}{b-a} = -j$, waarbij $j = \eu^{2\iu\pi/3}$ voldoet
aan $j^2 + j + 1 = 0$; of ontbind $a + jb + j^2c$ en $a + j^2 b + jc$.}
\end{exercise}

\begin{exercise}[$\star\star\star$]\label{exo:b1:complex:11}
Bereken voor $n \in \N^*$ het product $P = \prod_{k=1}^{n-1}
\bigl(1 - \omega^k\bigr)$, waarbij $\omega = \eu^{2\iu\pi/n}$.
\emph{Aanwijzing: $X^n - 1 = \prod_{k=0}^{n-1} (X - \omega^k)$; deel
door $X - 1$ en evalueer in $X = 1$.} Leid daaruit
$\prod_{k=1}^{n-1} \sin\frac{k\pi}{n} = \dfrac{n}{2^{n-1}}$ af.
\end{exercise}

\begin{exercise}[$\star\star$]\label{exo:b1:complex:12}
Zij $n \geq 2$. Los de vergelijking $(z + 1)^n = (z - 1)^n$ op in $\C$:
toon aan dat ze precies $n - 1$ oplossingen heeft, alle zuiver
imaginair, namelijk
\[
z_k = -\iu\,\frac{\cos\frac{k\pi}{n}}{\sin\frac{k\pi}{n}},
\qquad k = 1, \dots, n - 1 .
\]
\emph{Aanwijzing: $z = 1$ is geen oplossing, dus deel en gebruik de
\omterm{def:b1:complex:unity}{eenheidswortels}; pas daarna de halvehoekontbinding van
\cref{met:b1:complex:trig} toe.}
\end{exercise}

\section{Opgave: de stelling van Napoleon}

\begin{problem}\label{pb:b1:complex:1}
Zet op elke zijde van een \emph{willekeurige} driehoek naar buiten toe
een gelijkzijdige driehoek: de drie middelpunten daarvan vormen altijd
een gelijkzijdige driehoek. Dat is de stelling van
Napoleon\index{stelling van Napoleon} --- een \omterm{def:b1:logic:statement}{uitspraak} zonder
zichtbare reden om waar te zijn, die de algebra van $j = \eu^{2\iu\pi/3}$
in drie regels rekenwerk bewijst. Deze opgave bouwt het volledige
gereedschap op (draaiingen, directe gelijkvormigheden, het
$j$-criterium van \cref{exo:b1:complex:10}), bewijst de stelling van
Napoleon naar buiten en naar binnen toe, lokaliseert het ontaarde geval,
en sluit af met een tweede juweel van dezelfde school: de stelling van
van Schooten over gelijkzijdige driehoeken en de ongelijkheid van
Ptolemaeus. Overal worden punten van het vlak met hun affixen
vereenzelvigd.

\medskip
\noindent\textbf{Deel I --- Het getal $j$ en de draaiingen.}
\begin{enumerate}
  \item Bereken de algebraïsche vorm van $j$, en vervolgens $j^2$, $j^3$,
        $1 + j + j^2$, $\conj j$ en $j^{-1}$. Zet $1$, $j$, $j^2$ op een
        schets van de eenheidscirkel.
  \item Toon aan dat de draaiing met middelpunt $a$ over de hoek
        $\theta$ gegeven wordt door $r(z) = a + \eu^{\iu\theta}(z - a)$,
        en dat de samenstelling van twee draaiingen, over de hoeken
        $\theta$ en $\theta'$, een draaiing over $\theta + \theta'$ is
        wanneer $\theta + \theta' \notin 2\pi\Z$, en anders een
        translatie. (Gebruik \cref{prop:b1:complex:similitude}.)
  \item Zij $b \neq c$. Toon aan dat er precies twee punten $p$ zijn die
        $(b, c, p)$ gelijkzijdig maken, namelijk $p = b +
        \eu^{\pm\iu\pi/3}(c - b)$. Ga voor een \emph{directe} driehoek
        $(a, b, c)$ op het voorbeeld $a = 0$, $b = 1$, $c = \iu$ na dat
        de keuze $\eu^{-\iu\pi/3}$ die is welke aan de andere kant van
        de rechte $BC$ ligt dan $a$ --- de top \emph{naar buiten toe}.
  \item Volgens de oplossing van \cref{exo:b1:complex:10} is $(a, b, c)$
        direct gelijkzijdig precies wanneer $a + jb + j^2c = 0$. Bewijs
        de twee draaiingsidentiteiten
        \[
        b + jc + j^2a = j^2\,(a + jb + j^2c),
        \qquad
        c + ja + j^2b = j\,(a + jb + j^2c),
        \]
        en leid af dat het criterium invariant is onder cyclische
        \omterm{def:b1:counting:objects}{permutatie} van $(a, b, c)$.
  \item Toon aan dat, wanneer $a + jb + j^2c = 0$ en twee van de drie
        punten samenvallen, alle drie samenvallen. (Het criterium
        karakteriseert dus precies: directe gelijkzijdige driehoek, of
        één enkel punt.)
\end{enumerate}

\medskip
\noindent\textbf{Deel II --- Directe gelijkvormigheden.}
\begin{enumerate}[resume]
  \item Toon aan dat de \omterm{def:b1:logic:map}{afbeeldingen} $f(z) = az + b$ met $a \in \C^*$
        (de \emph{directe gelijkvormigheden}) stabiel zijn onder
        samenstelling en inversie: de samenstelling van twee ervan, en
        de inverse van elk ervan, is opnieuw van deze vorm. (In de taal
        van \cref{ch:b1:structures}: ze vormen een groep.)
  \item Toon aan dat er bij gegeven $z_1 \neq z_2$ en $w_1 \neq w_2$
        \emph{precies één} directe gelijkvormigheid $f$ bestaat met
        $f(z_1) = w_1$ en $f(z_2) = w_2$, en geef $a$ en $b$ expliciet.
  \item Toon aan dat een directe gelijkvormigheid $f(z) = az + b$ alle
        afstanden met $\abs a$ vermenigvuldigt en de \emph{vorm}
        $\frac{c - a'}{b' - a'}$ van elke driehoek $(a', b', c')$
        bewaart \emph{(teller en noemer worden allebei met $a$
        vermenigvuldigd)}. Leid af dat twee driehoeken direct
        gelijkvormig zijn precies wanneer hun vormen gelijk zijn.
  \item Bepaal volledig de directe gelijkvormigheid met $f(0) = 1$ en
        $f(1) = \iu$: geef $a$, $b$, het vaste punt, de factor en de
        hoek.
  \item Zij $\alpha + \beta = 1$. Toon aan dat het ``gewogen punt''
        $g(a', b') = \alpha a' + \beta b'$ met elke directe
        gelijkvormigheid verwisselt: $f(\alpha a' + \beta b') = \alpha
        f(a') + \beta f(b')$. Leid af dat zwaartepunten, middens en de
        napoleontische middelpunten hieronder alle door
        gelijkvormigheden worden meegevoerd --- de algebraïsche
        vergunning achter elk argument van het type ``zonder verlies van
        algemeenheid leggen we de omgeschreven cirkel op de
        eenheidscirkel''.
\end{enumerate}

\medskip
\noindent\textbf{Deel III --- De stelling van Napoleon.} Zij $(a, b, c)$
een directe driehoek. Zet op elke zijde de gelijkzijdige driehoek naar
buiten toe (vraag 3) en noem $n_a$, $n_b$, $n_c$ de middelpunten
(zwaartepunten) van de driehoeken op respectievelijk $[b, c]$, $[c, a]$
en $[a, b]$.
\begin{enumerate}[resume]
  \item Toon aan dat
        \[
        n_a = \alpha b + \beta c, \qquad
        n_b = \alpha c + \beta a, \qquad
        n_c = \alpha a + \beta b,
        \qquad\text{met}\quad
        \alpha = \frac{3 + \iu\sqrt3}{6},\ \beta = \conj\alpha .
        \]
  \item Ga $\alpha + \beta = 1$ na en leid af dat de driehoek $(n_a,
        n_b, n_c)$ \emph{hetzelfde zwaartepunt} heeft als $(a, b, c)$.
  \item Toon met de identiteiten van vraag 4 aan dat
        \[
        n_a + j\,n_b + j^2 n_c
        = (\alpha j^2 + \beta j)\,(a + jb + j^2c) .
        \]
  \item Bereken $\alpha j + \beta$ en besluit: $\alpha j^2 + \beta j =
        j(\alpha j + \beta) = 0$, zodat $n_a + j n_b + j^2 n_c = 0$ voor
        \emph{elke} driehoek: de buitenste napoleontische driehoek is
        direct gelijkzijdig (of één punt). Daarmee is de stelling van
        Napoleon bewezen.
  \item Zet de gelijkzijdige driehoeken in plaats daarvan \emph{naar
        binnen} toe (de keuze $\eu^{+\iu\pi/3}$ in vraag 3) en noem
        $m_a, m_b, m_c$ hun middelpunten. Toon aan dat $m_a = \beta b +
        \alpha c$ (en cyclisch), en bewijs vervolgens $m_a + j^2 m_b +
        j\,m_c = 0$: ook de binnenste napoleontische driehoek is
        gelijkzijdig, met de tegengestelde oriëntatie.
  \item Lokaliseer de ontaarding: toon aan dat $n_a = n_b = n_c$ precies
        gebeurt wanneer $a + j^2 b + jc = 0$, dat wil zeggen wanneer
        $(a, b, c)$ een \emph{indirecte} gelijkzijdige driehoek is.
        \emph{(Twee van de punten $n_a, n_b, n_c$ vallen samen precies
        wanneer alle drie samenvallen, wegens vraag 5; gebruik dan $n_a
        + j^2 n_b + j n_c = j(\alpha + \beta j)(a + j^2b + jc)$ en ga na
        dat $\alpha + \beta j \neq 0$.)}
  \item Voer de hele berekening uit op de driehoek $a = 0$, $b = 1$, $c
        = \iu$: geef $n_a, n_b, n_c$ exact, en ga door de drie
        gekwadrateerde zijden rechtstreeks te berekenen na dat de
        driehoek gelijkzijdig is met gekwadrateerde zijde
        $\frac{2 + \sqrt3}{3}$.
\end{enumerate}

\medskip
\noindent\textbf{Deel IV --- Ptolemaeus en van Schooten.}
\begin{enumerate}[resume]
  \item Bewijs de identiteit, geldig voor alle complexe $a, b, c, d$:
        \[
        (a - b)(c - d) + (a - d)(b - c) = (a - c)(b - d) .
        \]
  \item Leid daaruit de \emph{ongelijkheid van
        Ptolemaeus}\index{ongelijkheid van Ptolemaeus} af: voor elke
        vier punten $A, B, C, D$ geldt
        \[
        AC \cdot BD \;\leq\; AB \cdot CD + AD \cdot BC ,
        \]
        met gelijkheid dan en slechts dan als $(a-b)(c-d)$ en
        $(a-d)(b-c)$ op eenzelfde halfrechte vanuit $0$ liggen.
  \item Toon aan dat voor $\theta, \varphi \in \R$ geldt
        $\abs{\eu^{\iu\theta} - \eu^{\iu\varphi}} =
        2\,\abs{\sin\frac{\theta - \varphi}2}$.
  \item (Stelling van van Schooten\index{stelling van van Schooten})
        Zij $(a, b, c) = (1, j, j^2)$ --- wegens Deel II gaat daarmee
        geen algemeenheid verloren onder de directe gelijkzijdige
        driehoeken --- en zij $p = \eu^{\iu\theta}$ met $\theta \in
        \intoo{2\pi/3}{4\pi/3}$, een punt van de omgeschreven cirkel op
        de boog $BC$ die $A$ niet bevat. Bewijs
        \[
        PA = PB + PC .
        \]
        \emph{(Druk de drie afstanden uit met vraag 20 en gebruik een
        som-naar-productformule.)}
  \item Ga in $p = -1$ na dat dit precies het gelijkheidsgeval van de
        ongelijkheid van Ptolemaeus is voor de cyclische volgorde $A, B,
        P, C$: bereken $(a - b)(p - c)$ en $(a - c)(b - p)$ en
        controleer dat hun quotiënt een positief reëel getal is.
\end{enumerate}

\medskip
\noindent\textbf{Deel V --- Synthese.}
\begin{enumerate}[resume]
  \item Bereken de buitenste napoleontische driehoek van de
        gelijkzijdige driehoek $(1, j, j^2)$ zelf, en beschrijf het
        resultaat meetkundig.
  \item Waar precies gebruikte de opgave: (i) de vermenigvuldiging als
        draaiing; (ii) de groepsstructuur van de gelijkvormigheden en de
        invariante vorm; (iii) de halvehoekontbinding van
        \cref{met:b1:complex:trig}? Eén zin per onderdeel.
  \item Formuleer in een korte alinea de moraal van de opgave: wat
        levert het woordenboek tussen de vlakke meetkunde en de algebra
        van $\C$ op, wat kost het, en welke van de twee bewijsstappen
        --- de identiteit $\alpha j + \beta = 0$ of het klassieke
        plaatje --- \emph{verklaart} de stelling van Napoleon volgens
        jou beter? Noem één plaats waar dit woordenboek later in dit
        volume in matrixvorm terugkeert.
\end{enumerate}
\end{problem}
